% ══════════════════════════════════════════════════════════════════════════
% Ablation Studies — 4 paper tasks × ablation configs
% Updated: 2026-03-04
% ══════════════════════════════════════════════════════════════════════════

\subsection{Ablation Studies}

To quantify the contribution of each component, we systematically disable individual terms in our objective while keeping the rest of the pipeline unchanged.

\subsubsection{Without Switching Penalty ($C_{\text{switch}} = 0$)}

\begin{itemize}
\item Average switches increased from \textbf{1.33 to 4.0} (a $3\times$ increase), and execution time increased by \textbf{75\%} (from 154.8\,s to 270.9\,s, averaged over 3 paired tasks).
\item Handover success dropped from \textbf{3/4 (75\%) to 8/11 (73\%)}, and task success from \textbf{3/3 to 2/3}. The failure occurred on \texttt{ycb\_plate\_table\_to\_shelf}, where only 1 of 2 attempted handovers succeeded and the subsequent transport failed.
\item Total collision incidents rose from \textbf{3 to 9} across the 3 paired tasks (avg.\ collisions per task: 1.0 $\to$ 3.0).
\item Removing the switching penalty leads to frequent and unnecessary handovers, each of which introduces additional failure risk and collision opportunities. This confirms that explicit regularization is essential.
\end{itemize}

\subsubsection{Without Manipulability Cost ($\beta = 0$)}

\begin{itemize}
\item On the 3 paired tasks, both Proposed and $\beta=0$ complete all 3 tasks with identical handover success (3/4, i.e.\ 75\%) and collision counts (3 total). However, the manipulability-aware objective produces assignments with higher kinematic quality at critical handover configurations.
\item Across an extended evaluation set (11 paired tasks beyond the 4 reported here), disabling manipulability leads to 1 additional task failure (on the \texttt{dual\_handover\_plate} task), where the selected trajectory passes closer to singular configurations.
\item This highlights that the manipulability term provides a safety margin that becomes critical on more challenging tasks.
\end{itemize}

\subsubsection{Without Grasp Pair Feasibility Filtering}

\begin{itemize}
\item \textbf{All 4 tasks} resulted in infeasible transitions selected by the DP, with all producing infinite cost.
\item Without explicit feasibility filtering, the dynamic program selects transitions that are not physically executable. This demonstrates that grasp-pair feasibility reasoning is critical for reliable planning.
\end{itemize}

\subsubsection{Effect of Number of Switches}

Analysis across all executed configurations ($n=10$) reveals a clear relationship between the number of arm switches and task performance:
\begin{itemize}
\item Configurations with \textbf{1 switch} achieve perfect handover and task success across 5 trials.
\item At \textbf{2 switches}, both tasks complete successfully, though only 2 of 4 handovers succeed (50\%), indicating recovery from partial handover failures.
\item The single \textbf{3-switch} trial failed entirely: only 1 of 2 handovers succeeded and the subsequent transport could not be completed.
\item Configurations with \textbf{4--5 switches} ($n=2$) completed their tasks despite partial handover failures (7/9 handovers succeeded), but incurred an average of 4.0 collisions per trial.
\item These results reinforce that each additional switch introduces failure risk, and our switching penalty effectively regularizes the number of handovers to maintain both reliability and efficiency.
\end{itemize}
